\section{Population likelihoods and strategic interaction}
\label{app:population-proof}
For a fixed list, independence multiplies individual zero-arrival
probabilities $\exp(-\mu_{A,i})$, yielding
\begin{equation}
 P_0^{\rm fixed}=\exp(-\sum_i\mu_{A,i}).\label{eq:fixed}
\end{equation}
For finite origin intensity $\bar N_A$, condition on a Poisson number
$K$ of origins and independent marks drawn from the normalized intensity.
With $m=\mathbb E\exp(-\mu_A)$,
\[
 P_0=\sum_{k=0}^\infty e^{-\bar N_A}
 \frac{\bar N_A^k}{k!}m^k
 =\exp[-\bar N_A(1-m)],
\]
proving \cref{eq:origins}. Restriction to increasing finite-intensity
regions and continuity of probabilities proves the corresponding
sigma-finite formulation whenever needed, allowing a zero limit. This
is the generating-functional identity of \citet{LastPenrose2017}.

For \cref{prop:uniformity}, monotonicity of $1-e^{-u}$ implies
$\mathbb E(1-e^{-\mu_A})\geq\pi_{\rm out}(1-e^{-\mu_{\rm out}})$.
The hypothesis $P_0\geq P_*$ bounds this expectation by
$\mathcal L_*/\bar N_A$, proving the first inequality. Its inward contribution is
non-negative and bounded by the same total, giving the leakage
inequality. Two constant marks evaluate the expectation as
$f_{\rm in}+\pi_{\rm out}(f_{\rm out}-f_{\rm in})$.
If $\mathcal L_*/\bar N_A<f_{\rm in}$ no mixture works; if it is at least
$f_{\rm out}$ every mixture works. Otherwise the exact upper fraction
is $(\mathcal L_*/\bar N_A-f_{\rm in})/(f_{\rm out}-f_{\rm in})$.
At equal means the fraction cancels and only the common mean matters.
Since $0\leq1-e^{-\mu_A}\leq1$, $P_0\geq e^{-\bar N_A}$ proves the
universal rarity sufficient condition. A lower mean bound alone cannot
improve it: take $\bar N_A=100$, $P_*=.05$, and lower bound
$\mu_{\rm out}=.01$, but actual means twenty for every origin.
The proposed weaker rarity expression is below $\mathcal L_*$, while the actual
null probability is near $e^{-100}$. If $\bar N_A=0$, the null probability
is one without restrictions. If $\mu_{\rm out}=0$, division in
\cref{eq:uniformity-bound} is invalid and no bound follows from that
zero lower bound. Fractions zero and one and equality $P_0=P_*$ are
included by the weak inequalities. Inward leakage averages the full
function of the random mark, not that function at its mean.

\subsection{Proof of the preemption proposition}
\label{app:preemption-proof}
Fix rivals in \cref{eq:preemption-game}. The own-action objective is
strictly concave, so its unique maximizer satisfies exactly the three
conditions in \cref{eq:preemption-KKT}. Simultaneous satisfaction is
the definition of Nash equilibrium. The continuous function
$\Phi(x)=\sum_i f(x_i)+\iota\sum_{i<k}x_ix_k$ attains a maximum on
the compact cube. A unilateral change changes $\Phi$ by precisely the
change in that player's payoff, so such a maximum is an equilibrium.
Substituting all-zero and all-cap profiles proves their conditions;
substituting a common interior action gives the stated scalar equation.
This characterizes asymmetric profiles too through the coordinate
conditions; it does not assert that every equilibrium is symmetric.

For binary choices, the payoff difference between cap and zero is
$\bar X[d_{\rm bin}+\iota X_{-i}]$. Its sign proves the two endpoint
criteria. With one maintained seed and all others initially zero, the
remaining players initially see difference proportional to
$d_{\rm bin}+\iota\bar X$. If positive, each first updating player
adopts, and the incentive never decreases as others adopt. Under a fair
asynchronous update schedule, all adopt in finitely many individual
updates. If non-positive, none adopts with inward tie-breaking, so the
seed does not trigger a cascade. This proves the claimed iff under the
specified update and tie rules.

For the continuous counterexample, best response is
$\max\{0,(2x_{-i}-1)/10\}$ on $[0,1]$. At an equilibrium both actions
are at most $0.1$, which forces both responses to zero. Hence the sole
equilibrium is zero despite the positive origin derivative against a
maximal rival. The binary five-player counterexample follows by direct
substitution. These calculations refute an unconditional one-defector
cascade; they do not exclude cascades in games satisfying the stated
strict condition. The thermal identity in \cref{eq:radiator} is algebra,
not an additional behavioral theorem.
